\documentclass[10pt,a4paper]{amsart}
\usepackage[margin=19mm]{geometry}
\usepackage{amsmath,amssymb,mathtools}
\usepackage[scr=rsfs,cal=euler]{mathalfa}
\usepackage{xfrac}
\usepackage[unicode,hidelinks,bookmarks=false]{hyperref}
\newcommand{\ie}{i.e.\ }
\newcommand{\cf}{cf.\ }
\newcommand{\resp}{respectively\ }
\newcommand{\Subsection}{\subsection}
\providecommand{\nobreakdash}{\nobreak\hbox{-}}
\setlength{\emergencystretch}{2em}
\multlinegap0pt
\usepackage{amssymb}
\newtheorem{proposition}{Proposition}
\newcommand{\PP}{\mathbb{P}}
\newcommand{\cO}{\mathcal{O}}
\title{Polarized varieties without arithmetically Cohen--Macaulay bundles and section rings without graded maximal Cohen--Macaulay modules in characteristic zero}
\author{Cristian Anghel}
\date{Source: arXiv:2609.15589v1 (CC BY 4.0)}
\begin{document}
\maketitle

\noindent\emph{Source and licence.} Polarized varieties without arithmetically Cohen--Macaulay bundles and section rings without graded maximal Cohen--Macaulay modules in characteristic zero. Version arXiv:2609.15589v1 (CC BY 4.0), distributed by arXiv under the Creative Commons Attribution 4.0 International licence, http://creativecommons.org/licenses/by/4.0/, as recorded in the arXiv licence metadata for this exact version and shown on the abstract page https://arxiv.org/abs/2609.15589v1. Full licence notice: \texttt{licenses/arxiv-2609.15589-CC-BY-4.0.txt}.\par\smallskip

\noindent\textit{Editorial scope.} Counted unit: the article's proposition prop:bpf (source0.tex lines 287--289). The article's complete original proof of it is delivered with it (source0.tex lines 291--304, one \texttt{proof} environment of 14 lines). No mathematics is written, repaired or re-derived: the statement and the proof are the article's own lines, in the article's own notation, and no retained line is substituted or re-flowed.  Retained uncounted as the source-local context the proof needs: lines 76--76.  Nothing was omitted from any retained span, so the package carries no omission record.  No retained line contains a citation, so the wrapper carries no bibliography and no bibliography entry.  No retained line contains a figure environment, an image-inclusion command or drawing code, so no image is delivered and the package declares no asset dependency.  Editorial formatting only, in addition: an AMS \texttt{amsart} preamble that restates the article's own declarations and macros used by the retained lines, namely the environments \texttt{proposition} on separate counters, together with the standard packages those lines need.  The retained environments print in this document as Proposition 1 in order of appearance, whereas the article prints the same items with the numbering of its own full document.  The complete native source of the article as distributed in the arXiv e-print for this exact version is bundled unchanged as \texttt{sources/210430/source0.tex}.
\section{Statements}\label{sec:statements}

\begin{proposition}\label{prop:bpf}
For every $n\ge2$ the linear system $|H_n|$ is base-point-free and $H_n\cdot C>0$ for every irreducible curve $C\subset Y_n$; hence $\varphi_{H_n}$ is finite and $H_n$ is ample.
\end{proposition}

\begin{proof}
(a) $A_n=r^*\cO_{X_n}(1)$ is base-point-free.

(b) The coordinate $x_i$ of $\PP^{11}$ restricts to a section of $\cO_{X_n}(1)$ with $n\operatorname{div}(x_i)=\operatorname{div}(x_i^n)=p^*\ell_i$; hence $r^*x_i\in H^0(Y_n,\cO(A_n))$ satisfies $n\operatorname{div}(r^*x_i)=f^*\beta^*\ell_i=f^*\bigl(\tilde L_i+\sum_{p\in\ell_i}E_p\bigr)=n\bigl(D_i+\sum_{p\in\ell_i}E'_p\bigr)$, i.e.
\[
\operatorname{div}(r^*x_i)=D_i+\sum_{p\in\ell_i}E'_p .
\]

(c) Let $T$ be one of the four triangles of the Hesse pencil. Its three lines contain the nine flexes, each flex lying on exactly one of them (the vertices of the four triangles are the twelve double points). By (b), $\Phi_T:=\prod_{i\in T}r^*x_i\in H^0(Y_n,\cO(3A_n))$ has divisor $\sum_{i\in T}D_i+E_n$; hence $\Phi_T=\Psi_T\cdot z$, where $z$ is the canonical section of $\cO(E_n)$ and $\Psi_T\in H^0(Y_n,\cO(3A_n-E_n))$ has divisor $\sum_{i\in T}D_i$.

(d) Two lines $\ell_i\in T$ and $\ell_j\in T'$ from different triangles meet at a flex: $\ell_i$ carries three flexes, each lying on exactly one line of $T'$, and these three lines are distinct (a line of $T'$ through two flexes of $\ell_i$ would be $\ell_i$); so every line of $T'$ meets $\ell_i$ at one of its flexes. Since distinct lines through a blown-up point have disjoint strict transforms, $\tilde L_i\cap\tilde L_j=\emptyset$ and $D_i\cap D_j=f^{-1}(\tilde L_i\cap\tilde L_j)=\emptyset$. Hence for $T\ne T'$ the sections $\Psi_T$ and $\Psi_{T'}$ have disjoint zero loci, and $|3A_n-E_n|$ is base-point-free.

(e) $H_n=A_n+(3A_n-E_n)$ is base-point-free, as a sum of base-point-free divisors. It is ample: on $Z$ one has $H_Z=4L-\sum_pE_p=-K_Z+L$, where $-K_Z=3L-\sum_pE_p$, the class of the Hesse pencil, is nef because the pencil has no base points on $Z$, and $L$ is nef; so $H_Z\cdot E_p=1$, $H_Z\cdot C\ge L\cdot C>0$ for every other irreducible curve $C$, and $H_Z^2=7>0$, whence $H_Z$ is ample by the Nakai--Moishezon criterion, and so is $H_n=\tfrac1nf^*H_Z$, $f$ being finite. An ample base-point-free divisor defines a finite morphism (Section~\ref{sec:statements}).
\end{proof}

\end{document}
